%-------------------------------------------------------------------------------
% Mẫu CV Chuẩn Harvard (Bản Tiếng Việt)
% Tối ưu cho hệ thống ATS và bố cục chuẩn 1 trang
% Ứng viên: HUỲNH CÔNG LUYỆN
% Vị trí mục tiêu: Junior AI Engineer | Junior RAG Engineer | Junior ML/AI Engineer
%                  Backend AI Engineer | AI Chatbot Engineer
%-------------------------------------------------------------------------------

\documentclass[10pt,letterpaper]{article}

% --- Engine Detection & Hỗ trợ Tiếng Việt / Unicode ---
\usepackage{iftex}
\ifPDFTeX
  \usepackage[utf8]{vietnam}
\else
  \usepackage{fontspec}
\fi

% --- Cấu hình lề trang (Đo lường chính xác cho 1 trang duy nhất) ---
\usepackage[
  letterpaper,
  top=0.33in,
  bottom=0.33in,
  left=0.42in,
  right=0.42in
]{geometry}

% --- Gói định dạng văn bản ---
\usepackage{titlesec}
\usepackage{enumitem}
\usepackage{xcolor}
\usepackage{hyperref}
\usepackage{tabularx}
\usepackage[verbose=silent]{microtype}
\linespread{1.04}

% --- Bảng màu ---
\definecolor{primary}{RGB}{0, 0, 0}
\definecolor{secondary}{RGB}{45, 45, 45}
\definecolor{linkcolor}{RGB}{18, 65, 145}

% --- Cấu hình Hyperref ---
\hypersetup{
  colorlinks=true,
  linkcolor=linkcolor,
  filecolor=linkcolor,
  urlcolor=linkcolor,
  pdfauthor={Huỳnh Công Luyện},
  pdftitle={Huỳnh Công Luyện - CV Tiếng Việt},
  pdfsubject={Kỹ sư AI / RAG CV},
  pdfkeywords={RAG, Vector Search, LLM, Rust, Python, Qdrant, Lượng tử hóa}
}

% --- Bỏ thụt lề và số trang ---
\pagestyle{empty}
\raggedbottom
\setlength{\tabcolsep}{0in}

% --- Định dạng tiêu đề mục ---
\titleformat{\section}{
  \vspace{-4pt}\raggedright\large\bfseries\color{primary}
}{}{0em}{}[\color{black}\titlerule \vspace{2pt}]

% --- Định dạng danh sách ---
\newcommand{\resumeItemListStart}{\begin{itemize}[leftmargin=1.15em, label={\small$\bullet$}, itemsep=1.1pt, parsep=0pt, topsep=0.6pt, partopsep=0pt]}
\newcommand{\resumeItemListEnd}{\end{itemize}\vspace{-1pt}}

\newcommand{\resumeItem}[1]{\item \small{#1}}

\newcommand{\resumeSubheading}[4]{
  \item[]
  \begin{tabular*}{\textwidth}[t]{l@{\extracolsep{\fill}}r}
    \textbf{#1} & \textbf{\small #2} \\
    \textit{\small #3} & \textit{\small #4} \\
  \end{tabular*}\vspace{1pt}
}

\begin{document}

%===============================================================================
% THÔNG TIN CÁ NHÂN / LIÊN HỆ
%===============================================================================
\begin{center}
  {\Large \textbf{Huỳnh Công Luyện}} \\ \vspace{2pt}
  {\large \textbf{\color{secondary}KỸ SƯ AI / RAG}} \quad{\small$|$\quad \textbf{Workflow RAG $\bullet$ Tích hợp LLM $\bullet$ Dịch vụ Backend}} \\ \vspace{2.5pt}
  \small{
    Email: \href{mailto:huynhluyen839@gmail.com}{huynhluyen839@gmail.com} \quad$\bullet$\quad
    Điện thoại: (+84) 926072506 \quad$\bullet$\quad
    Địa điểm: TP. Hồ Chí Minh, Việt Nam
  } \\ \vspace{1.5pt}
  \small{
    GitHub: \href{https://github.com/neshaki091}{github.com/neshaki091} \quad$\bullet$\quad
    Portfolio: \href{https://Neshaki091.vercel.app}{Neshaki091.vercel.app} \quad$\bullet$\quad
    Trạng thái: Sẵn sàng làm việc toàn thời gian
  }
\end{center}
\vspace{-3pt}

%===============================================================================
% TÓM TẮT CHUYÊN MÔN
%===============================================================================
\section{Tóm tắt Chuyên môn}
\small{
Sinh viên IT mới tốt nghiệp (GPA 3.34/4.0) với kinh nghiệm thực tế xây dựng hệ thống RAG/chatbot, tích hợp LLM API, đánh giá chất lượng prompt/output và triển khai dịch vụ backend. Tác giả \textbf{ARQ-RAG}, engine truy xuất vector giảm $18,8\times$ RAM hoạt động so với FAISS trên 5M vector (Khóa luận: 9.0/10), và dẫn dắt dự án \textbf{EduChatbot}, nơi em trực tiếp thiết kế workflow RAG và triển khai backend trên GCP Cloud Functions.
}
\vspace{3.5pt}

%===============================================================================
% HỌC VẤN
%===============================================================================
\section{Học vấn}
\begin{itemize}[leftmargin=0pt, label={}]
  \resumeSubheading
    {Trường Đại học Giao thông vận tải TP. Hồ Chí Minh (UTH)}{TP. Hồ Chí Minh, Việt Nam}
    {Cử nhân Công nghệ Thông tin \textnormal{(Đã hoàn thành chương trình; Chờ cấp bằng)}}{2022 -- 2026}
    \resumeItemListStart
      \resumeItem{\textbf{GPA:} 3.34 / 4.0 \quad$\bullet$\quad \textbf{Khóa luận:} ARQ-RAG --- Điểm bảo vệ: 9.0 / 10.0}
    \resumeItemListEnd
\end{itemize}
\vspace{0pt}

%===============================================================================
% KỸ NĂNG CỐT LÕI
%===============================================================================
\section{Kỹ năng Cốt lõi}
\begin{itemize}[leftmargin=0pt, label={}, itemsep=1.0pt]
  \item \small{\textbf{AI \& RAG:} Truy xuất ngữ nghĩa, Embeddings, Prompt Engineering, Đánh giá output, Viết lại truy vấn, Đối chiếu trích dẫn chống ảo giác, LangChain, QLoRA, Ragas}
  \item \small{\textbf{Tìm kiếm Vector:} Qdrant, Supabase pgvector, Chỉ mục IVF, Lượng tử hóa Vector, Ánh xạ bộ nhớ (\texttt{mmap}), Đánh giá truy xuất}
  \item \small{\textbf{Backend \& Cloud:} JavaScript (Node.js, Express), Python (FastAPI), Java (Spring Boot), Docker Compose, Nginx, MongoDB, Redis, GCP Cloud Functions, Git}
  \item \small{\textbf{Ngoại ngữ:} Tiếng Anh: A2 (đọc hiểu tài liệu kỹ thuật và papers) $\bullet$ Tiếng Việt: Bản ngữ}
\end{itemize}
\vspace{2pt}

%===============================================================================
% DỰ ÁN & NGHIÊN CỨU TIÊU BIỂU
%===============================================================================
\section{Dự án \& Nghiên cứu Tiêu biểu}
\begin{itemize}[leftmargin=0pt, label={}]

  % --- DỰ ÁN 1: ARQ-RAG ---
  \item[]
  \begin{tabular*}{\textwidth}[t]{l@{\extracolsep{\fill}}r}
    \textbf{\large ARQ-RAG --- Engine Truy xuất Vector Thích ứng} & \textbf{\small 2026} \\
    \textit{\small Tác giả chính \quad$|$\quad Thiết kế thuật toán \& Hiện thực Rust bằng AI \quad$|$\quad \href{https://arq-rag-frontend.vercel.app/}{Live Demo}} & \textit{\small Khóa luận: \textbf{9.0 / 10.0}} \\
  \end{tabular*}\vspace{0pt}
  \resumeItemListStart
    \resumeItem{\textbf{Vai trò:} Thiết kế kiến trúc hệ thống (chỉ mục IVF, lượng tử hóa phần dư, định tuyến cụm thích ứng). Sử dụng AI agent sinh mã Rust hiệu năng cao (AVX2 SIMD, \texttt{mmap}); tự mình review và xác thực toàn bộ mã nguồn.}
    \resumeItem{\textbf{RAM:} Giảm RAM hoạt động từ 1.931 MB (FAISS) xuống \textbf{102,6 MB} ($18,8\times$) trên 5 triệu vector Wikipedia DPR ($D$=768, CPU đơn luồng) bằng cách chỉ lưu mã 4-bit trên RAM; vector gốc truy xuất từ đĩa qua zero-copy \texttt{mmap}.}
    \resumeItem{\textbf{Tốc độ \& Recall:} Đạt \textbf{4.522 QPS} (87,2\% Recall@16 trước rerank) và \textbf{99,2\% Recall@16} sau khi Float32 rerank tại 3.450 QPS. Lập chỉ mục 500K vector trong 20,2s (nhanh $1,93\times$ so với FAISS IVFPQ).}
    \resumeItem{\textbf{Học được:} Lý thuyết lượng tử hóa vector, phân tích đánh đổi bộ nhớ (RAM vs. đĩa vs. cache), và xây dựng benchmark có khả năng tái tạo. \textit{Bài báo đã submit AMI 2026 [Đang bình duyệt]}.}
  \resumeItemListEnd
  \vspace{3.5pt}

  % --- DỰ ÁN 2: EduChatbot ---
  \item[]
  \begin{tabular*}{\textwidth}[t]{l@{\extracolsep{\fill}}r}
    \textbf{\large EduChatbot --- Trợ lý Hỏi đáp Tài liệu RAG Tiếng Việt} & \textbf{\small Thg 06/2026} \\
    \textit{\small Trưởng nhóm (6 thành viên) \quad$|$\quad Python, FastAPI, Java 17, Spring Boot, Qdrant, Llama 3.1 8B \quad$|$\quad \href{https://github.com/Neshaki091/Educhatbot}{GitHub}} & \textit{\small Đồ án Môn học Java} \\
  \end{tabular*}\vspace{0pt}
  \resumeItemListStart
    \resumeItem{\textbf{Pipeline RAG \& Tích hợp:} Thiết kế workflow RAG toàn diện (parsing/phân đoạn tiếng Việt, embedding 384 chiều, truy xuất vector Qdrant, đối chiếu trích dẫn chống ảo giác); xây cầu nối FastAPI tích hợp với Spring Boot.}
    \resumeItem{\textbf{Đánh giá, Cloud \& Trưởng nhóm:} Điều phối nhóm 6 người (theo dõi sprint, review mã nguồn) và deploy GCP Cloud Functions kết hợp Supabase + Qdrant Cloud; đánh giá A/B Ragas chứng minh RAG đạt độ trung thực cao hơn (0,303 vs. 0,292) và độ trễ thấp hơn (2,00s vs. 3,75s) so với fine-tuning.}
  \resumeItemListEnd
  \vspace{3.5pt}

  % --- DỰ ÁN 3: Waterbase ---
  \item[]
  \begin{tabular*}{\textwidth}[t]{l@{\extracolsep{\fill}}r}
    \textbf{\large Waterbase --- Nền tảng BaaS Đa người dùng Tự host} & \textbf{\small Thg 11/2025} \\
    \textit{\small Lập trình viên duy nhất \quad$|$\quad Node.js, Express, MongoDB, Redis, Docker Compose, Nginx \quad$|$\quad \href{https://tttn-project-waterbase-mini-clone-f.vercel.app/}{Live Demo}} & \textit{\small Đồ án Thực tập Tốt nghiệp} \\
  \end{tabular*}\vspace{0pt}
  \resumeItemListStart
    \resumeItem{\textbf{Kiến trúc Backend \& Triển khai:} Tự thiết kế và điều phối nền tảng microservices gồm 7 services (Auth, App, WaterDB, Realtime, Storage, Rule Engine, Analytics) qua Nginx Gateway, vận hành thực tế ổn định trên VPS 1 GB RAM / 1 vCPU.}
    \resumeItem{\textbf{Multi-Tenancy \& Phân quyền:} Triển khai mô hình cô lập dữ liệu Database-per-Application triệt tiêu rủi ro rò rỉ dữ liệu, xác thực JWT kèm xoay vòng refresh token, và Rule Engine (RBAC) phân quyền khai báo chi tiết cấp collection.}
  \resumeItemListEnd
  \vspace{3.5pt}

  % --- DỰ ÁN 4: Trợ lý AI Nội bộ ---
  \item[]
  \begin{tabular*}{\textwidth}[t]{l@{\extracolsep{\fill}}r}
    \textbf{\large Trợ lý AI Nội bộ --- Lotus Hackathon VNG 2026} & \textbf{\small Thg 03/2026} \\
    \textit{\small Lập trình viên Backend (đội 2 người) \quad$|$\quad Node.js, Express, Supabase pgvector, Gemini API, n8n \quad$|$\quad \href{https://devpost.com/software/internal-ai-assistance}{Video Demo}} & \textit{\small Tập đoàn VNG} \\
  \end{tabular*}\vspace{0pt}
  \resumeItemListStart
    \resumeItem{\textbf{Backend AI Thần tốc (36h):} Xây dựng backend Node.js/Express hoàn chỉnh trong \textbf{36 giờ}, tích hợp RAG trên Supabase \texttt{pgvector} qua Gemini, cô lập dữ liệu multi-tenant bằng \texttt{clientId}, và router dự phòng ý định 3 tầng (parser $\rightarrow$ Gemini $\rightarrow$ từ khóa) kết hợp n8n.}
  \resumeItemListEnd

\end{itemize}

\end{document}
